\scanpage{8}
\mspage{7}
\begin{center}4 - 27 -\end{center}
\margin{n\textsuperscript{o} 392}

Soit maintenant $L_1$ un second sous-module libre de $M$, de rang $n$. Puisque $S^{-1}L=S^{-1}L_1$, il existe $s\in S$ tel que $sL_1\subset L$ ; on peut donc se borner à prouver que si $L_1\subset L_2$ sont deux sous-modules libres de $M$, on a $c(\chi(M/L_1))=c(\chi(M/L_2))$. Or, on a $\chi(M/L_1)=\chi(M/L_2)+\chi(L_2/L_1)$, et il résulte du n\textsuperscript{o} $6$, cor. de la prop. $14$, que $\chi(L_2/L_1)$ est un diviseur principal, donc $c(\chi(L_2/L_1))=0$. C. Q. F. D.

L'élément $c(\chi(M/L))$ sera noté $-c(M)$ dans ce qui suit ; nous dirons que $c(M)$ est la \uline{classe de diviseurs attachée à} $M$.

PROPOSITION 16.- (i) \uline{Soit} $0\longrightarrow M_1\xrightarrow{f}M_2\xrightarrow{g}M_3\longrightarrow0$ \uline{une suite exacte de} $A$-\uline{modules de type fini}. \uline{On a} :
\[
 c(M_2)=c(M_1)+c(M_3).
\]

(ii) \uline{S'il existe un pseudo-isomorphisme de} $M_1$ \uline{dans} $M_2$, \uline{on a} $c(M_1)=c(M_2)$.

(iii) \uline{Si} $T$ \uline{est un module de torsion}, \uline{on a} $c(T)=-c(\chi(T))$.

(iv) \uline{Si} $\underline a$ $\ne0$ \uline{est un idéal fractionnaire de} $K$, \uline{on a} $c(\underline a)=c(\operatorname{div}(\underline a))$.

(v) \uline{Si} $L$ \uline{est un} $A$-\uline{module libre}, \uline{on a} $c(L)=0$.

Pour prouver (i), considérons un sous-module libre $L_1$ (resp. $L_3$), de $M_1$, (resp. $M_3$), tel que $M_1/L_1$, (resp. $M_3/L_3$), soit un module de torsion. Puisque $L_3$ est libre et $g$ surjectif, il existe dans $g^{-1}(L_3)$ un supplémentaire libre $L_{23}$ de $\operatorname{Ker}(g)$, isomorphe à $L_3$, (\uline{Alg.}, chap. II, 3e éd., \S\ $1$, n\textsuperscript{o} $11$, prop. $21$) ; mais $\operatorname{Ker}(g)=f(M_1)$ contient $f(L_1)=L_{12}$ qui est libre puisque $f$ est injectif. La somme $L_2=L_{12}+L_{23}$ est directe, et $L_2$ est donc un sous-module \uline{libre} de $M_2$. On a en outre le diagramme commutatif :
\diagram{\begin{tikzcd}
0\arrow[r]&L_1\arrow[r]\arrow[d]&L_2\arrow[r]\arrow[d]&L_3\arrow[r]\arrow[d]&0\\
0\arrow[r]&M_1\arrow[r,"f"']&M_2\arrow[r,"g"']&M_3\arrow[r]&0
\end{tikzcd}}
